\section{A joint gluing inequality}
\label{sec:gluing}

The geometric construction will reach a set of possible relay vertices and
then need to continue to a target. The vertex reached is selected by the
same edge configuration that must supply the continuation, so its
connection probability cannot simply be treated as an independent factor.
We prove an inequality that controls this dependence uniformly in the
number of relays.

Connections in this section are witnessed by finite open paths, including
paths of length zero.  For an undirected graph $G$ with independent edge
states, write $C^G(x)$ for the vertex cluster of $x$, and put
$C^G(S)=\bigcup_{x\in S}C^G(x)$.  Superscripts are omitted when the graph
is clear.  Vertex deletion retains the original probabilities on all
remaining edges.  Parallel edges are allowed; loops can be discarded.

\begin{theorem}[Joint gluing]
\label{thm:joint-gluing}
Let $G$ have countably many vertices and edges, with independent edge-open
probabilities in $[0,1]$.  For a vertex $o$, a nonempty set $A$ of vertices,
and any set $T$ of vertices,
\begin{equation}
\label{eq:joint-gluing}
 \mathbb P(o\leftrightarrow A,\ o\leftrightarrow T)
 \ge \mathbb P(o\leftrightarrow A)
          \inf_{a\in A}\mathbb P(a\leftrightarrow T).
\end{equation}
For finite $A$, the infimum is a minimum.
\end{theorem}

Taking $T=\{b\}$ and enlarging the event on the left proves the
multiplicative gluing inequality of Kozma and Nitzan,
Conjecture~1 of \cite{KozmaNitzan2024}. The finite statement also follows
from the first-contact comparison (GEN) in \cite[Section~3.4]{Leder2026},
as recorded in \cite{ConnectionInequalities2026}. We give a direct proof,
then pass to countable graphs.

For the finite proof, with $o\notin A$, we will construct numbers
$\alpha_a\ge0$, $a\in A$, whose sum is $\mathbb P(o\leftrightarrow A)$
and for which
\[
 \mathbb P(o\leftrightarrow A,\ o\leftrightarrow T)
 \ge \sum_{a\in A}\alpha_a\mathbb P(a\leftrightarrow T).
\]
Bounding the sum by its total mass times the least target probability
will prove the theorem. The weights come from the inverse of a triangular
matrix of conditional cluster probabilities. A positivity property of
its cluster columns will give both the required signs in the inverse
and the displayed comparison. Related coefficient
questions appear in \cite[Section~5.2]{KozmaNitzan2024}, where the system
uses a symmetric matrix of joint connection probabilities instead.

\subsection{Conditional cluster columns}

Throughout the finite-graph proof assume $p_e<1$ for every edge; zero
probabilities are permitted.  If $C(S)=w$ has positive probability, its
occurrence is determined by edges having an endpoint in $w$: all vertices
of $w$ must connect to $S$ within $w$, and all edges from $w$ to its
complement must be closed.  Thus edges induced outside $w$ retain their
independent original laws.  The same observation applies inside a graph
obtained by vertex deletion.  For $S=\varnothing$, its cluster is empty.

Choose distinct distinguished vertices $x_1,\ldots,x_N$.
They need not exhaust the vertices of $G$; all other vertices and their
edges remain in the graph.  For $1\le k\le N$, put
\[
 S_k=\{x_1,\ldots,x_{k-1}\},\qquad
 E_k=\{C(x_k)\cap S_k=\varnothing\},\qquad d_k=\mathbb P(E_k)>0.
\]
The positivity follows by closing every edge incident to $x_k$.
Let $\mu_k$ be the law of $C(x_k)$ conditional on $E_k$, and set
\[
 \mathcal D_k=\{U\subseteq V(G):x_k\in U,\ U\cap S_k=\varnothing\}.
\]
For $F:\mathcal D_k\to[0,\infty)$, define vectors in $\mathbb R^N$ by
\[
 v_k^F(i)=\mathbb E_{\mu_k}
              [\mathbf1_{\{x_i\in C(x_k)\}}F(C(x_k))],
 \qquad h_k=v_k^1.
\]
The matrix $H=(h_1\ \cdots\ h_N)$ is unit lower triangular.

The vector $v_k^F-(\mathbb E_{\mu_k}F)h_k$ records the conditional
covariances of $F(C(x_k))$ with the distinguished-vertex indicators.
For increasing $F$, we will show that this vector is a nonnegative
combination of later columns.
The following deletion estimate supplies the nonnegative terms in that
combination.

\subsection{The deletion estimate and column positivity}

\begin{lemma}[Deleted-set residual]
\label{lem:deleted-residual}
Fix a finite product graph $G$, a vertex $s$, a vertex set $S$ not
containing $s$, and an increasing function $F$ on sets that contain $s$
and avoid $S$.  For $w\supseteq S$ with $s\notin w$, define
\[
 c(w)=\mathbb E^{G\setminus w}F(C(s)).
\]
For $U\subseteq V(G)\setminus(S\cup\{s\})$, sample
$W_U=C^{G\setminus U}(S)$ and define
\[
 R_F(U)=\mathbb E\left[
  \mathbf1_{\{s\notin W_U\}}
       \bigl(c(W_U)-c(W_U\cup U)\bigr)\right],
\]
where the integrand is zero when $s\in W_U$.  Then $R_F$ is nonnegative
and increasing under inclusion of $U$.
\end{lemma}

\begin{proof}
Sample all edges of $G$ once and use their restrictions for every graph
below.  Conditional on an exact value $W_U=w$, the edges induced outside
$w$ have product law, including those involving the previously deleted
vertices $U$.  Consequently $R_F(U)$ is the expectation of
\[
 D_U=\mathbf1_{\{s\notin W_U\}}
 \left[F\bigl(C^{G\setminus W_U}(s)\bigr)
       -F\bigl(C^{G\setminus(W_U\cup U)}(s)\bigr)\right]
\]
in this common configuration.  This quantity is nonnegative.

If $U\subseteq U'$, then $W_{U'}\subseteq W_U$.  On
$\{s\notin W_U\}$, the first cluster in $D_U$ grows when $U$ is replaced
by $U'$.  The second cluster is exactly $C^{G\setminus U}(s)$: in
$G\setminus U$ the cluster of $s$ is disjoint from the union of the
$S$-clusters.  The corresponding second cluster for $U'$ is
$C^{G\setminus U'}(s)$ and can only shrink.  Increasingness of $F$
therefore gives $D_{U'}\ge D_U$ on this event.  On its complement,
$D_U=0$ and $D_{U'}\ge0$.  Taking expectations proves the claim.
\end{proof}

\begin{lemma}[Cluster-column positivity]
\label{lem:cluster-column}
For every nonnegative increasing $F$ on $\mathcal D_k$,
\begin{equation}
\label{eq:cluster-column}
 v_k^F-(\mathbb E_{\mu_k}F)h_k
       \in\left\{\sum_{l>k}a_lh_l:a_l\ge0\right\}.
\end{equation}
In particular $H^{-1}v_k^F$ has nonnegative coordinates.
\end{lemma}

\begin{proof}
We induct downward in $k$. At a fixed index, exposing the earlier
clusters will split $v_k^F$ into an averaged version of itself and
nonnegative contributions from later columns. Repeating the averaging
will replace $F$ by its mean and leave exactly the difference in
\eqref{eq:cluster-column}.

For $k=N$, both vectors in the difference have only their $N$th
coordinate nonzero, so the difference vanishes. Fix $k<N$ and assume
the claim for all larger indices.  Write
$s=x_k$, $S=S_k$, $W=C(S)$, and
$c(w)=\mathbb E^{G\setminus w}F(C(s))$ for $w\supseteq S$ avoiding $s$.
The function $c$ is nonincreasing in $w$ by restriction coupling.

Condition on $W=w$ and $s\notin w$, and use the independent product
law outside $w$.  If $x_i\notin w$, partition its cluster by its earliest
distinguished vertex.  This is $x_k$ exactly when $x_i\in C(s)$;
write $L_{il}^w$ for the case where it is $x_l$, $l>k$.
Conditional on the exact cluster $U=C^{G\setminus w}(x_l)$ in the latter
case, the expected value of $F(C(s))$ is $c(w\cup U)$.
Partitioning the expectation $c(w)$ in this way yields
\begin{equation}
\label{eq:cluster-partition}
\begin{split}
 \mathbb E^{G\setminus w}
    [\mathbf1_{\{x_i\in C(s)\}}F(C(s))]
 &=c(w)\mathbb P^{G\setminus w}(x_i\in C(s))\\
 &\quad+\sum_{\substack{l>k\\x_l\notin w}}\mathbb E^{G\setminus w}
  [\mathbf1_{L_{il}^w}\{c(w)-c(w\cup C(x_l))\}].
\end{split}
\end{equation}
If $x_i\in w$, every term is zero instead.

We first average the leading term in \eqref{eq:cluster-partition}.
Given $C(s)=U$ under $E_k$, the earlier-root cluster $W$ has the law of
$W_U=C^{G\setminus U}(S)$, by the product law outside the exposed
cluster. Define
\[
 (PF)(U)=\mathbb E[c(W_U)],\qquad U\in\mathcal D_k.
\]
This function is nonnegative and increasing, since $W_U$ shrinks with
$U$ and $c$ decreases with its argument. Conditioning first on $W$ and
then on $C(s)$ shows that the averaged leading term is
\[
 \mathbb E[\mathbf1_{\{x_i\in C(s)\}}c(W)\mid E_k]
 =\mathbb E_{\mu_k}[\mathbf1_{\{x_i\in C(s)\}}(PF)(C(s))]
 =v_k^{PF}(i).
\]

The remaining terms give later columns. More precisely, averaging
\eqref{eq:cluster-partition} conditional on $E_k$ gives
\begin{equation}
\label{eq:column-decomposition}
 v_k^F=v_k^{PF}+\sum_{l>k}v_l^{F_l},\qquad
 F_l(U)=\frac{d_l}{d_k}R_F(U),\quad U\in\mathcal D_l.
\end{equation}
Here the residual uses the fixed roots $s,S$ of this induction step.
To check the normalization, on $E_k$ the event $L_{il}^W$ is exactly
$E_l\cap\{x_i\in C(x_l)\}$.  Given $C(x_l)=U$ under $E_l$, the
$S$-cluster is $W_U$ with product law in $G\setminus U$, and the remaining
condition $E_k$ is $s\notin W_U$.  Division by $d_k$ and conditioning
on $E_l$ give precisely the factor $d_l/d_k$.
Lemma~\ref{lem:deleted-residual} shows that every $F_l$ is nonnegative
and increasing.  The induction hypothesis therefore places every term
$v_l^{F_l}$ in the cone generated by $h_j$ with $j>k$.

Thus one application of $P$ leaves a remainder in the cone of later
columns. To iterate this decomposition, we identify the limit of $P^jF$.
The operator $P$ first samples $W$ conditional on $C(s)$ and then samples
$C(s)$ conditional on $W$. Both are conditional laws of the joint
distribution given $E_k$, so $\mu_k$ is stationary for this two-step
Markov chain. These alternating cluster updates are the method of
van den Berg, H\"aggstr\"om and
Kahn~\cite[Section~2.1]{VanDenBergHaggstromKahn2006}.
Every transition puts mass at least
\[
 a=\prod_{e\ni s}(1-p_e)>0
\]
on $\{s\}$.  Subtracting this common mass from two transition averages
shows $\operatorname{osc}(Pg)\le(1-a)\operatorname{osc}(g)$ for every
function on the finite state space.  Stationarity then implies
$P^jF\to\mathbb E_{\mu_k}F$ uniformly.

Iterate \eqref{eq:column-decomposition}.  Each $P^jF$ remains
nonnegative and increasing, so $v_k^F-v_k^{P^jF}$ belongs to the cone
generated by the later columns.  This cone is closed: on coordinates
$k+1,\ldots,N$ its generators form a unit lower-triangular basis.
Taking the limit proves \eqref{eq:cluster-column}.  Finally
$H^{-1}h_l$ is the $l$th coordinate vector, proving the last assertion.
\end{proof}

\subsection{Nonnegative relay weights}

Apply Lemma~\ref{lem:cluster-column} to
\[
 F(U)=\mathbf1_{\{U\cap\{x_{k+1},\ldots,x_N\}\ne\varnothing\}}.
\]
Its mean $m_k$ under $\mu_k$ is less than one, since $x_k$ can be
isolated.  The centered vector is $(1-m_k)(h_k-e_k)$, where $e_k$
is the $k$th coordinate vector.  Hence there are numbers
$\Gamma_{lk}\ge0$, $l>k$, such that
\[
 h_k=e_k+\sum_{l>k}\Gamma_{lk}h_l.
\]
Let $\Gamma$ be zero on and above the diagonal.  In matrix form,
$H=I+H\Gamma$, and therefore
\begin{equation}
\label{eq:cluster-inverse}
 H^{-1}=I-\Gamma.
\end{equation}
Moreover, partitioning the cluster of $x_i$ by its first distinguished
vertex gives
\begin{equation}
\label{eq:cluster-row-mass}
 Hd=\mathbf1,\qquad (I-\Gamma)\mathbf1=d,
 \qquad d=(d_1,\ldots,d_N)^{\mathsf T}.
\end{equation}

\begin{proof}[Proof of Theorem~\ref{thm:joint-gluing}]
First suppose the graph is finite and all $p_e<1$.
If $o\in A$, the left side of \eqref{eq:joint-gluing} equals
$\mathbb P(o\leftrightarrow T)$ and the inequality is immediate.
Otherwise take the distinguished list to consist of the vertices of
$A$ followed by $x_N=o$, and set
\[
 F(U)=\mathbf1_{\{U\cap A\ne\varnothing,\ U\cap T\ne\varnothing\}},
 \qquad y_i=\mathbb E F(C(x_i)).
\]
This is one nonnegative increasing function, restricted to each
$\mathcal D_k$ as needed.  Partitioning by the first distinguished
vertex gives $y=\sum_k d_kv_k^F$, so
Lemma~\ref{lem:cluster-column} implies $H^{-1}y\ge0$ coordinatewise.
The last row of \eqref{eq:cluster-inverse} yields
\[
 y_N\ge\sum_{i<N}\Gamma_{Ni}y_i.
\]
For $i<N$, length-zero paths give $y_i=\mathbb P(x_i\leftrightarrow T)$,
whereas $y_N$ is the joint probability in the theorem.  By
\eqref{eq:cluster-row-mass},
\[
 \sum_{i<N}\Gamma_{Ni}=1-d_N=\mathbb P(o\leftrightarrow A).
\]
Thus the promised relay weights are $\alpha_{x_i}=\Gamma_{Ni}$ for $i<N$.
They are nonnegative, proving the finite-graph inequality.
This argument also covers $T=\varnothing$, when $F=0$.

All event probabilities on a finite graph are polynomials in the edge
probabilities.  Approximating each probability equal to one from below
therefore includes deterministic open edges as well as closed ones.

For a countable graph and finite $A$, exhaust the graph by finite
subgraphs containing $o$ and $A$, and intersect $T$ with their vertex
sets.  Every connection has a finite witness, so the individual and
joint connection events increase to the corresponding events in $G$.
The minimum over the finite set $A$ also converges.  This proves
\eqref{eq:joint-gluing} for finite $A$.  Finally exhaust a countable
nonempty $A$ by finite nonempty sets $A_j$.  The two events involving
$A_j$ increase to those involving $A$, while
$\min_{a\in A_j}\mathbb P(a\leftrightarrow T)$ decreases to the infimum
over $A$.  Taking limits proves the result.
\end{proof}

\begin{corollary}[Failure after reaching a relay]
\label{cor:gluing-failure}
In the setting of Theorem~\ref{thm:joint-gluing}, let $0\le\lambda\le1$.
If $\mathbb P(a\leftrightarrow T)\ge1-\lambda$ for every $a\in A$, then
\[
 \mathbb P(o\leftrightarrow A,\ o\not\leftrightarrow T)
 \le\lambda\,\mathbb P(o\leftrightarrow A)\le\lambda.
\]
The same conclusion holds for $A=\varnothing$.
\end{corollary}

\begin{proof}
Subtract \eqref{eq:joint-gluing} from $\mathbb P(o\leftrightarrow A)$.
For an empty relay set the event on the left is empty.
\end{proof}

The theorem and corollary may be applied after fixing any collection of
edge states, provided the remaining edges still have product law.
Relay and target sets must then be fixed under that conditioning.
In particular, a later adaptive choice of contacts is not an additional
conditioning to be inserted into their connection probabilities.

We apply this inequality in the corridor construction of
Section~\ref{sec:corridor-exploration}. The remaining sections develop
the geometric argument for graphs satisfying the polynomial bound along
an unbounded sequence of radii, beginning with the quotient action that
will supply the corridor coordinates.
